\scanpage{171}
\uline{Groupe fondamental} : (Revêtements infinis)

$S$ préschéma. Pour tout groupe discret $G$, on définit
\[
H^1_{\unsure{\mathcal C|\mathscr C}}(S,G)=\varinjlim_{\substack{T/S\in\mathcal C\\ \deleted{\unsure{T\ \text{fid. plat de type fini sur }S}}\\ \deleted{\unsure{\text{de type qct sur }S|\text{de type qct vers }S}}}}H^1(T/S,G)
\]
où
\[
H^1(T/S,G)=H^1(K_{T/S},G)=H^1(\pi_0(K_{T/S}),G)
\]
et où $\mathcal C$ est une catégorie de $S$-préschémas.

% slanted
\margin{élucider ici les questions d'effectivité des données de descente\unsure{.....} \unsure{Se ramener|Ça ramène} à la descente des morphismes étales de type fini\unsure{,....}}

Les cas les plus importants sont les suivants

\diagram{\(\begin{array}{l}
\left\{\begin{array}{l}
\text{(i)}\quad \mathcal C_1\qquad S\text{-préschémas fid plats quasi-compacts, }\deleted{\unsure{\text{localemt}|\text{localement}}}.\\
\left.\begin{array}{lll}
\text{(ii)}&\mathcal C_2&\rule{4em}{0.3pt}\quad\text{de type fini}\\
\text{(iii)}&\mathcal C_3&\rule{4em}{0.3pt}\quad\text{quasi-finis}
\end{array}\right\}
\end{array}\right.\\[3pt]
\left.\begin{array}{lll}
\text{(iv)}&\mathcal C_4&\rule{4em}{0.3pt}\quad\text{finis par morceaux}\\
\text{(v)}&\mathcal C_5&\rule{4em}{0.3pt}\deleted{\text{étales}}\quad\text{finis étales par morceaux}
\end{array}\right\}
\end{array}\)}

On a donc $H^1_{\mathcal C_i}(\deleted{\unsure{T,|P}}S,G)\hookrightarrow H^1_{\mathcal C_j}(S,G)$ si $i\geq j$
application injective canonique. \uline{Si $G$ est fini},
alors tous les hom sont bijectifs, et $H^1(S,G)$
est l'ensemble des classes (à isom près) des
revêtements étales galoisiens de groupe $G$, donc $\simeq$
$\operatorname{Hom}(\pi_1(S,a),G)/\operatorname{int}(G)$. Notre intérêt ici est pour
le cas $G$ infini. Notons en passant
\[
H^1_{\mathcal C_3}\xrightarrow{\sim}H^1_{\mathcal C_2}
\]
% archive sheet number: 170
